% !TeX root = ../../main.tex
\section{状态依赖误差的执行示例与验证方法}
\subsection{一个两步收敛的精确算术示例}
为说明 CG 状态之间的依赖，考虑二维对称正定系统
\begin{equation}
 A=\begin{bmatrix}1&0\\0&4\end{bmatrix},\qquad
 b=\begin{bmatrix}1\\1\end{bmatrix},\qquad x_0=0.
\end{equation}
下述计算使用精确分数运算，是解释算法的构造示例，用于分析 CG 的状态依赖关系。初始残差与搜索方向均为 $(1,1)^{\mathsf T}$。第一轮 $q_0=Ap_0=(1,4)^{\mathsf T}$，因此
\begin{equation}
 \alpha_0=\frac{r_0^{\mathsf T}r_0}{p_0^{\mathsf T}q_0}=\frac25,
 \quad x_1=\begin{bmatrix}2/5\\2/5\end{bmatrix},
 \quad r_1=\begin{bmatrix}3/5\\-3/5\end{bmatrix}.
\end{equation}
进一步得到 $\beta_0=9/25$，搜索方向为 $p_1=(24/25,-6/25)^{\mathsf T}$。第二轮步长为 $\alpha_1=5/8$，于是 $x_2=(1,1/4)^{\mathsf T}$，残差为零。该例展示精确算术下二维 CG 的有限步性质，实际浮点系统中的行为还受到舍入和实现细节影响。

\begin{table}[htbp]\centering\caption{二维示例中的精确 CG 状态}
\begin{tabular}{ccccc}\toprule
轮次&搜索方向&矩阵向量积&步长&更新后残差\\\midrule
0&$(1,1)$&$(1,4)$&$2/5$&$(3/5,-3/5)$\\
1&$(24/25,-6/25)$&$(24/25,-24/25)$&$5/8$&$(0,0)$\\\bottomrule
\end{tabular}\end{table}

这个过程中的后续方向由前一轮残差决定，而前一轮残差又依赖矩阵向量积和步长。因此，混合精度不是对独立 SpMV 任务分别作选择；同一轮的误差会改变下一轮输入，构成闭环控制需要考虑的状态演化。

\subsection{单次乘积扰动如何改变后续状态}
现在只在第一轮把乘积替换为 $\widetilde q_0=(1,9/2)^{\mathsf T}$，即人为加入 $e_0=(0,1/2)^{\mathsf T}$。该扰动是用于分析的构造值，与特定浮点格式的实际舍入误差无关。

新的步长为 $\widetilde\alpha_0=4/11$，更新解为 $(4/11,4/11)^{\mathsf T}$，递推残差为 $(7/11,-7/11)^{\mathsf T}$。但由原矩阵重新计算的真实残差是 $(7/11,-5/11)^{\mathsf T}$，二者间隙为
\begin{equation}
 g_1=\begin{bmatrix}0\\2/11\end{bmatrix}
 =\widetilde\alpha_0e_0.
\end{equation}
这与前述残差间隙关系一致。递推残差的范数并不等于真实残差范数，新的 $\beta$和搜索方向也与精确过程不同。即使下一轮恢复高精度，也是在已经改变的状态上继续执行。

该例有两个启示。第一，检查单轮乘积的相对误差不足以直接预测完整求解时间；它还通过步长、搜索方向和残差影响后续轮次。第二，终止时应明确检查的是递推量还是真实残差。学习策略可以利用状态相关性提高效率，但不能替代独立的数值质量验证。

\subsection{局部误差与矩阵结构的联系}
对混合精度乘积可写为 $\widetilde q=(A+\Delta A)p+\delta$，其中 $\Delta A$概括数值表示变化，$\delta$概括乘加及归约舍入。该分解仅是解释性记号，具体误差界需要根据实际内核和数据类型推导。相同位宽在不同 Tile 上产生的绝对扰动取决于矩阵值、搜索方向幅值、行内抵消和归约长度。

若某 Tile 中的搜索方向分量很小，该区域的乘积绝对贡献可能较低；但在存在强烈抵消的行中，即使每个局部贡献不大，其变化仍可能明显改变最后的行结果。因此，“当前值小”只能作为可能的精度容忍特征，不能作为普遍的安全降精度准则。状态特征与训练环境共同决定策略能否利用这些关系。

同样，矩阵条件数能够提示问题的数值敏感性，却不能单独决定某轮某 Tile 应使用何种精度。求解状态随迭代变化，局部向量分布也随之变化。本文采用搜索方向片段提供动态信息，但其状态表征没有包含所有可能影响收敛的因素，这也是迁移和鲁棒性需要实验验证的原因。

\subsection{精度比例的加权统计方法}
若一半 Tile 采用低精度，并不表示一半非零元素或一半计算时间采用低精度。Tile 之间的非零数和访问成本可能不同。可以分别定义按 Tile 数、按逻辑非零数以及按实际读取字节数加权的精度比例：三者对应决策分布、算术工作量和访存需求，回答的问题不同。

例如两个 Tile 分别包含100 和900 个非零元素，若仅第一个 Tile 采用 BF16，则按 Tile 数统计低精度比例为50\%，按非零数统计则为10\%。这里同样是构造示例。若存在补齐，物理槽位比例还会进一步变化。分析实验的精度比例时，应保持论文所采用的统计口径，不把它自动解释成等比例的带宽节省。

\begin{table}[htbp]\centering\caption{混合精度统计的不同加权口径}
\begin{tabularx}{\textwidth}{p{2.8cm}XX}\toprule
统计方式&反映的信息&解释限制\\\midrule
按动作或 Tile 数&策略选择的空间分布&不反映每个区域的大小\\
按逻辑非零数&参与运算的数据规模&忽略填充与元数据\\
按物理槽位数&实际数值流长度&仍未计入全部事务开销\\
按测得执行时间&关键路径上的成本构成&重叠执行时不宜直接相加\\\bottomrule
\end{tabularx}\end{table}

\subsection{学习策略与规则策略的公平比较}
比较学习策略和固定规则时，应保持它们可见的状态与使用的内核基础尽可能一致。若规则实现通过 CPU 决定动作，而学习策略完全在 GPU 运行，时间差中同时包含决策质量和执行位置差异。为了隔离决策能力，可补充同设备执行的规则；为了评价完整系统，则保留端到端实现并说明不同路径。

规则阈值也需要在独立验证集合上选择。如果学习策略经过大量训练，而规则仅采用一个未经校准的阈值，比较会混合调参预算差异。反过来，若规则使用测试结果逐矩阵挑选最佳参数，也会引入有利偏差。合理协议应提前规定搜索范围、训练或验证预算，以及最终测试时是否允许调整。

实验提供了规则基线的成功数量和时间对照，可用于说明特定规则在测试条件下的局限。该比较的结论限定于所采用的规则及其参数配置。学习策略的优势仍应通过相同执行基础上的比较、失败案例和训练成本共同评估，并与可验证的数值保护机制互补。

\subsection{运行时优化的收益归属}
统一内核、向量操作融合和 CUDA Graph 均可降低执行开销，即便所有矩阵值采用 FP64 也可能受益。因此，本章首先用 EvoCG-FP64 作为相同执行基础的参照，再比较动态精度。相对通用库的完整系统加速比用于展示系统效果，相同执行基础上的加速比用于分析精度控制的增量价值，两者分别刻画系统整体收益与精度控制的增量贡献。

进一步区分时，可以保持矩阵布局不变，分别启用策略推理、动作复用和低精度数值流。但消融配置可能改变执行重叠和缓存行为，各模块时间差不一定具有严格可加性。若将每个独立收益相乘，也可能重复计算相互依赖的改善。因此，评估采用列出完整配置组合及总时间，并用时间分解辅助解释。

CUDA Graph 还要求重复执行结构与数据依赖满足捕获和重放条件。动态动作若只改变设备数组中的精度标记，可以在固定图结构中被内核读取；若每轮动作决定完全不同的内核数量，则可能增加图更新成本。EvoCG 的统一执行路径使动态精度尽量表现为数据变化，这是算法粒度与运行时机制协同的具体体现。

\subsection{从矩阵专属策略到任务族复用}
同一矩阵配合不同右端项时，搜索方向轨迹仍可能改变；不同初值、停止阈值和预条件器也会改变状态分布。因而“矩阵相同”是策略复用的重要条件，却未必充分。部署前应明确训练配置与实际求解配置的差异，并观察成功率、迭代数和真实残差，而不只记录模型推理是否正常执行。

任务族复用可以降低训练摊销压力，但应把训练矩阵、验证矩阵和最终测试矩阵分开。若先观察所有测试矩阵再挑选一个迁移效果最好的源策略，其报告结果会包含选择偏差。已有跨矩阵实验揭示了可复用机会及失败边界，下一步需要更严格的任务族划分与多种子验证，而不能仅增加训练样本后假定泛化问题已经解决。

对求解失败的回退同样需要明确状态。可以从原初值使用 FP64 重新求解，也可以在当前近似解上重新计算真实残差并重启；两者具有不同成本和数值行为。这些回退方案的完整性能与数值验证有待后续研究。它们的共同要求是保留原问题和独立质量检查，使性能优化始终服务于相同的数值求解目标。
